% GENERATED FILE. Edit canonical lessons, then run npm run generate:books.
% canonical-source-hash: fnv1a64:6c4fe74f5a7605dd
% canonical-lessons: TE-C415-vidya, TE-C415-jnanam, TE-C415-adhyayanam, TE-C415-tarkincu, TE-C415-bodhincu

\chapter{Learning and Teaching}
\label{ch:te-learning-and-teaching}
\hlchaptermodality{415}

\begin{chapteropening}
\textbf{By the end of this chapter:} \emph{I can talk about education, knowledge, study and teaching.}

Complete the last lesson of chapter 415: I can talk about education, knowledge, study and teaching.
\end{chapteropening}

\section[\texorpdfstring{\te{విద్య} (vidya)}{విద్య (vidya)}]{\te{విద్య} (vidya) — education; learning}

\label{lesson:TE-C415-vidya}



\begin{quote}
Before the new one: say the Telugu for a scholarship, then the Telugu for to question.
\end{quote}

\subsection*{You'll want to know: \te{విద్య}}
\textbf{\te{విద్య}} — \emph{vidya} — \textquotedblleft{}education; learning\textquotedblright{}.

\begin{quote}
\te{ప్రతి} \te{పిల్లవాడికి} \te{విద్య} \te{అవసరం}.
\end{quote}

\begin{quote}
\emph{prati pillavāḍiki vidya avasaraṁ.}
\end{quote}

\begin{quote}
Every child needs education.
\end{quote}

\te{విద్యార్థి}, a student, and \te{విశ్వవిద్యాలయం}, a university, which you know, are built on it.

\begin{cousinweb}
\textbf{\te{విద్య}} (\emph{vidya}) is borrowed from Sanskrit \emph{vidyā}, \textquotedblleft{}knowledge\textquotedblright{} (CDIAL 11741), from \emph{vid}, \textquotedblleft{}to know\textquotedblright{}; Telugu shortens the final long a.
\end{cousinweb}

\begin{grammarlens}[title={What you've built}]
One more word for giving reasons and opinions.
\end{grammarlens}

\subsection*{Your turn}
\begin{itemize}
\raggedright
  \item \emph{Say it:} \emph{vidya}
  \item \emph{Say it:} \emph{vidya}, once more
  \item \emph{Say it:} \emph{prati pillavāḍiki vidya avasaraṁ.} — Every child needs education.
  \item \emph{Recall:} say the Telugu for a scholarship, then the Telugu for to question, then say \emph{vidya} again
\end{itemize}

\begin{tcolorbox}[breakable,colback=teal!4,colframe=teal!35!black,title={Before you move on}]
What does \te{విద్య} mean? (\textquotedblleft{}Education; learning\textquotedblright{}.) Say it once more.
\end{tcolorbox}
\section[\texorpdfstring{\te{జ్ఞానం} (jñānaṁ)}{జ్ఞానం (jñānaṁ)}]{\te{జ్ఞానం} (jñānaṁ) — knowledge; wisdom}

\label{lesson:TE-C415-jnanam}



\begin{quote}
Before the new one: say the Telugu for to question, then the Telugu for education.
\end{quote}

\subsection*{You'll want to know: \te{జ్ఞానం}}
\textbf{\te{జ్ఞానం}} — \emph{jñānaṁ} — \textquotedblleft{}knowledge; wisdom\textquotedblright{}.

\begin{quote}
\te{పుస్తకాలు} \te{చదివితే} \te{జ్ఞానం} \te{పెరుగుతుంది}.
\end{quote}

\begin{quote}
\emph{pustakālu cadivitē jñānaṁ perugutundi.}
\end{quote}

\begin{quote}
If you read books, your knowledge grows.
\end{quote}

\te{జ్ఞాని}, a wise person.

\begin{cousinweb}
\textbf{\te{జ్ఞానం}} (\emph{jñānaṁ}) is borrowed from Sanskrit \emph{jñāna}, \textquotedblleft{}knowledge\textquotedblright{} (CDIAL 5281), from \emph{jñā}, \textquotedblleft{}to know\textquotedblright{}.
\end{cousinweb}

\begin{grammarlens}[title={What you've built}]
One more word for giving reasons and opinions.
\end{grammarlens}

\subsection*{Your turn}
\begin{itemize}
\raggedright
  \item \emph{Say it:} \emph{jñānaṁ}
  \item \emph{Say it:} \emph{jñānaṁ}, once more
  \item \emph{Say it:} \emph{pustakālu cadivitē jñānaṁ perugutundi.} — If you read books, your knowledge grows.
  \item \emph{Recall:} say the Telugu for to question, then the Telugu for education, then say \emph{jñānaṁ} again
\end{itemize}

\begin{tcolorbox}[breakable,colback=teal!4,colframe=teal!35!black,title={Before you move on}]
What does \te{జ్ఞానం} mean? (\textquotedblleft{}Knowledge; wisdom\textquotedblright{}.) Say it once more.
\end{tcolorbox}
\section[\texorpdfstring{\te{అధ్యయనం} (adhyayanaṁ)}{అధ్యయనం (adhyayanaṁ)}]{\te{అధ్యయనం} (adhyayanaṁ) — study; a study}

\label{lesson:TE-C415-adhyayanam}



\begin{quote}
Before the new one: say the Telugu for education, then the Telugu for knowledge.
\end{quote}

\subsection*{You'll want to know: \te{అధ్యయనం}}
\textbf{\te{అధ్యయనం}} — \emph{adhyayanaṁ} — \textquotedblleft{}study; a study\textquotedblright{}.

\begin{quote}
\te{ఆ} \te{అధ్యయనంలో} \te{కొత్త} \te{విషయాలు} \te{తెలిశాయి}.
\end{quote}

\begin{quote}
\emph{ā adhyayanaṁlō kotta viṣayālu teliśāyi.}
\end{quote}

\begin{quote}
New facts came to light in that study.
\end{quote}

\te{అధ్యయనం} \te{చేయు}, to study a subject closely.

\begin{cousinweb}
\textbf{\te{అధ్యయనం}} (\emph{adhyayanaṁ}) is borrowed from Sanskrit \emph{adhyayana}, \textquotedblleft{}reading, studying\textquotedblright{} (Monier-Williams).
\end{cousinweb}

\begin{grammarlens}[title={What you've built}]
One more word for giving reasons and opinions.
\end{grammarlens}

\subsection*{Your turn}
\begin{itemize}
\raggedright
  \item \emph{Say it:} \emph{adhyayanaṁ}
  \item \emph{Say it:} \emph{adhyayanaṁ}, once more
  \item \emph{Say it:} \emph{ā adhyayanaṁlō kotta viṣayālu teliśāyi.} — New facts came to light in that study.
  \item \emph{Recall:} say the Telugu for education, then the Telugu for knowledge, then say \emph{adhyayanaṁ} again
\end{itemize}

\begin{tcolorbox}[breakable,colback=teal!4,colframe=teal!35!black,title={Before you move on}]
What does \te{అధ్యయనం} mean? (\textquotedblleft{}Study; a study\textquotedblright{}.) Say it once more.
\end{tcolorbox}
\section[\texorpdfstring{\te{తర్కించు} (tarkiñcu)}{తర్కించు (tarkiñcu)}]{\te{తర్కించు} (tarkiñcu) — to debate; to reason}

\label{lesson:TE-C415-tarkincu}



\begin{quote}
Before the new one: say the Telugu for knowledge, then the Telugu for study.
\end{quote}

\subsection*{You'll want to know: \te{తర్కించు}}
\textbf{\te{తర్కించు}} — \emph{tarkiñcu} — \textquotedblleft{}to debate; to reason\textquotedblright{}.

\begin{quote}
\te{వాళ్ళు} \te{రాత్రంతా} \te{ఆ} \te{ప్రశ్న} \te{గురించి} \te{తర్కించారు}.
\end{quote}

\begin{quote}
\emph{vāḷḷu rātrantā ā praśna guriñci tarkiñcāru.}
\end{quote}

\begin{quote}
They debated that question all night.
\end{quote}

Past (I …): \te{తర్కించాను}, \emph{tarkiñcānu}.

The verb of \te{తర్కం}, logic.

\begin{cousinweb}
\textbf{\te{తర్కించు}} (\emph{tarkiñcu}) is built on Sanskrit \emph{tarka}, \textquotedblleft{}inquiry\textquotedblright{} (CDIAL 5714), \textquotedblleft{}reasoning\textquotedblright{} (Monier-Williams), with the Telugu verb ending -\te{ఇంచు}; the noun is \te{తర్కం}.
\end{cousinweb}

\begin{grammarlens}[title={What you've built}]
One more word for giving reasons and opinions.
\end{grammarlens}

\subsection*{Your turn}
\begin{itemize}
\raggedright
  \item \emph{Say it:} \emph{tarkiñcu}
  \item \emph{Say it:} \emph{tarkiñcu}, once more
  \item \emph{Say it:} \emph{vāḷḷu rātrantā ā praśna guriñci tarkiñcāru.} — They debated that question all night.
  \item \emph{Recall:} say the Telugu for knowledge, then the Telugu for study, then say \emph{tarkiñcu} again
\end{itemize}

\begin{tcolorbox}[breakable,colback=teal!4,colframe=teal!35!black,title={Before you move on}]
What does \te{తర్కించు} mean? (\textquotedblleft{}To debate; to reason\textquotedblright{}.) Say it once more.
\end{tcolorbox}
\section[\texorpdfstring{\te{బోధించు} (bōdhiñcu)}{బోధించు (bōdhiñcu)}]{\te{బోధించు} (bōdhiñcu) — to instruct; to teach}

\label{lesson:TE-C415-bodhincu}



\begin{quote}
Before the new one: say the Telugu for study, then the Telugu for to debate.
\end{quote}

\subsection*{You'll want to know: \te{బోధించు}}
\textbf{\te{బోధించు}} — \emph{bōdhiñcu} — \textquotedblleft{}to instruct; to teach\textquotedblright{}.

\begin{quote}
\te{ఆమె} \te{పిల్లలకు} \te{గణితం} \te{బోధిస్తుంది}.
\end{quote}

\begin{quote}
\emph{āme pillalaku gaṇitaṁ bōdhistundi.}
\end{quote}

\begin{quote}
She teaches mathematics to children.
\end{quote}

Past (I …): \te{బోధించాను}, \emph{bōdhiñcānu}.

Formal beside \te{నేర్పు}, which you know; the person taught takes -\te{కి}.

\begin{cousinweb}
\textbf{\te{బోధించు}} (\emph{bōdhiñcu}) is built on Sanskrit \emph{bodhana}, \textquotedblleft{}causing to awake; enlightening, teaching\textquotedblright{} (Monier-Williams), with the Telugu verb ending -\te{ఇంచు}.
\end{cousinweb}

\begin{grammarlens}[title={What you've built}]
One more word for giving reasons and opinions.
\end{grammarlens}

\subsection*{Your turn}
\begin{itemize}
\raggedright
  \item \emph{Say it:} \emph{bōdhiñcu}
  \item \emph{Say it:} \emph{bōdhiñcu}, once more
  \item \emph{Say it:} \emph{āme pillalaku gaṇitaṁ bōdhistundi.} — She teaches mathematics to children.
  \item \emph{Recall:} say the Telugu for study, then the Telugu for to debate, then say \emph{bōdhiñcu} again
\end{itemize}

\begin{tcolorbox}[breakable,colback=teal!4,colframe=teal!35!black,title={Before you move on}]
What does \te{బోధించు} mean? (\textquotedblleft{}To instruct; to teach\textquotedblright{}.) Say it once more.
\end{tcolorbox}
